\par\Needspace{6\baselineskip}\medskip
\begin{center}\hypertarget{native-177-article130}{}\pdfbookmark[2]{130.}{article-130}130.\end{center}
Decem methodi inde ab art. 124 explicatae innituntur suppositioni, valores
approximatos distantiarum corporis coelestis a terra, aut positionis plani
orbitae, iam in potestate esse. Quoties quidem id agitur, ut dimensiones
orbitae, quarum valores approximati iam alicunde innotuerunt, puta per calculum
anteriorem observationibus aliis innixum, per observationes magis ab invicem
remotas corrigantur, postulatum illud nullis manifesto difficultatibus obnoxium
erit. Sed hinc nondum liquet, quonam modo calculum primum aggredi liceat,
ubi omnes orbitae dimensiones penitus adhuc incognitae sunt: hic vero
problematis nostri casus longe gravissimus atque difficillimus est, uti iam
ex problemate analogo in theoria cometarum praesumi potest, quod quamdiu
geometras torserit quotque tentaminibus irritis originem dederit satis constat.
Ut problema nostrum recte solutum censeri possit, manifesto conditionibus
sequentibus satisfieri oportet, siquidem solutio ad instar normae inde ab
art. 119 explicatae exhibetur: \emph{Primo} quantitates $x,y$ tali modo sunt
eligendae, ut valores ipsarum approximatos ex ipsa problematis natura petere
liceat, saltem, quamdiu corporis coelestis motus heliocentricus intra
observationes non nimis magnus est. \emph{Secundo} autem requiritur, ut variationibus
exiguis quantitatum $x,y$ variationes non nimis magnae in quantitatibus inde
derivandis respondeant, ne errores in illarum valoribus suppositis forte
commissi impediant, quominus has quoque pro approximatis habere liceat.
Denique \emph{tertio} postulamus, ut operationes, per quas a quantitatibus $x,y$
successive usque ad $X,Y$ progrediendum est, non nimis prolixae evadant.

Hae conditiones criterium subministrabunt, secundum quod de cuiusvis
methodi praestantia iudicium ferri poterit: adhuc evidentius quidem ea
applicationibus frequentibus se manifestabit. Methodus ea, quam exponere iam
accingimur, et quae quodammodo tamquam pars gravissima huius operis
consideranda est, illis conditionibus ita satisfacit, ut nihil amplius desiderandum
\SourceLeaf{178}relinquere videatur. Quam antequam in forma ad praxin commodissima explicare
aggrediamur, quasdam considerationes praeliminares praemittemus, aditumque
quasi ad illam, qui alias forsan obscurior minusque obvius videri possit,
illustrabimus atque aperiemus.
